\newpage
\section{ANEXOS}

Los anexos del trabajo se entregan como documentos independientes y no forman parte del cómputo de páginas de esta memoria. Cada anexo es autocontenido y la memoria los referencia donde corresponde para evitar duplicación de contenido.

\begin{table}[H]
\centering
\small
\caption{Anexos del Trabajo de Grado.}
\label{tab:anexos}
\begin{tabular}{p{1.5cm}p{6.0cm}p{7.5cm}}
\toprule
\textbf{Id.} & \textbf{Documento} & \textbf{Contenido principal} \\
\midrule
A & Plan de Gestión del Proyecto (SPMP) & Calendarización, presupuesto, riesgos, calidad, configuración y cierre. \\
B & Especificación de Requisitos (SRS) & Requerimientos funcionales y no funcionales, casos de uso, modelo de datos, criterios de aceptación. \\
C & Diseño (SDD) y Propuesta de TG & Vistas C4, secuencias UML, modelo on-chain/off-chain y justificación tecnológica. \\
\bottomrule
\end{tabular}
\end{table}

\textbf{Anexo A. Plan de Gestión del Proyecto (SPMP).} Contiene la vista general del proyecto, los supuestos y restricciones, los entregables, la calendarización ejecutada, el presupuesto, los riesgos materializados, las prácticas de calidad, la gestión de configuración y el cierre retrospectivo del proyecto. Incluye cinco diagramas BPMN del proceso de gestión y la lista completa de hitos cumplidos hasta el cierre del semestre 2026-1.

\textbf{Anexo B. Especificación de Requisitos de Software (SRS).} Contiene la definición detallada de los veinte requerimientos funcionales y los veinte no funcionales del prototipo organizados por módulo, los casos de uso por actor, el modelo de datos relacional, las interfaces del backend con Soroban RPC y Horizon, los criterios de aceptación y la trazabilidad entre requerimientos, operaciones del contrato y pruebas unitarias.

\textbf{Anexo C. Especificación del Diseño y Propuesta del Trabajo de Grado.} Documento que consolida la propuesta original del trabajo, el documento de diseño (SDD) con las vistas C4 (contexto, contenedores, componentes y despliegue), los diagramas de secuencia de los flujos críticos (compra primaria, reventa atómica y redención) y el modelo de datos completo on-chain/off-chain. Incluye el análisis de alternativas tecnológicas y la justificación de la selección de Stellar/Soroban frente a otras plataformas de contratos inteligentes.

\textbf{Material complementario.} Repositorio público del proyecto: \url{https://github.com/Tesis-Stellar/Secure-Ticket}.
